% Appendix C. Source: Appendix 4.B of Dissertation/Chapter4_Methodology.md.

\chapter{Assumptions Behind the Approval Wait}
\label{app:queue}
\lhead{Appendix C. \emph{Assumptions Behind the Approval Wait}}

Equation~(\ref{eq:pk}) is the Pollaczek-Khinchine formula for an M/G/1 queue. The formula gives the exact long-run mean waiting time when its assumptions hold. Low utilisation reduces the size of the predicted wait. Table~\ref{tab:mg1} sets out each assumption, what it means, and how it applies to this service.

{\small
\begin{xltabular}{\textwidth}{@{}P{0.17\textwidth}>{\hsize=0.90\hsize}L>{\hsize=1.10\hsize}LP{0.17\textwidth}@{}}
  \caption{The M/G/1 assumptions checked against this service}
  \label{tab:mg1}\\
  \toprule
  \textbf{Assumption} & \textbf{What it requires} & \textbf{Evidence for this service} & \textbf{Status} \\
  \midrule
  \endfirsthead
  \toprule
  \textbf{Assumption} & \textbf{What it requires} & \textbf{Evidence for this service} & \textbf{Status} \\
  \midrule
  \endhead
  \midrule
  \multicolumn{4}{@{}r@{}}{\footnotesize\emph{continued overleaf}}\\
  \endfoot
  \bottomrule
  \endlastfoot

    Poisson arrivals &
      Work reaches the Work reaches the approver independently and is
      dispersed over time rather than arriving in batches spread out, not in
      bunches &
      A trial's clock starts at HRA approval or site selection, both set by
      the sponsor and the regulator rather than by the trust. Trials therefore
      arrive at their own pace. The fixed committees do not batch the
      signature, for the reason given below &
      Met, subject to the batching rationale below \\
    \addlinespace
    General service time &
      Only the mean and the coefficient of variation are needed &
      Pharmacy has been asked to provide direct information on approval time.
      The formula itself does not require approval times to follow a normal,
      lognormal, or other specified distribution. &
      Met \\
    \addlinespace
    Single server &
      One approver, with work queueing behind whatever they are already doing &
      The trust describes one approving person for pharmacy set-up. Pharmacy
      adds (August 2026 return) that the sign-off can be delegated to the
      operational manager when most processes are already approved &
      Approximation, conservative: a second signer can only shorten the wait \\
    \addlinespace
    Work-conserving, non-preemptive &
      The approver is never idle with work waiting, and never abandons a
      part-finished approval &
      This is consistent with the described process. It also means that the
      mean wait does not depend on the order in which trials are approved. &
      Met \\
    \addlinespace
    Unbounded queue &
      The approver does not remain idle while approval work is waiting and
      does not abandon a partly completed approval &
      Trials wait for the signature. None are abandoned for want of one &
      Met \\
    \addlinespace
    Independent service times &
      The duration of one approval does not depend on the duration of other
      approvals or on the number waiting &
      Not observable in the available records. &
      Assumed, recorded as a limitation \\
    \addlinespace
    Steady state, $\rho < 1$ &
      On average, the approver can complete work faster than new approval
      requests arrive. &
      The commitment cycle resets monthly, so the service may not fully settle
      into a long-run average. &
      Assumed, recorded as a limitation 
\end{xltabular}}


The coefficient of variation, $\kappa$, is not identified by the available records. Equation~(\ref{eq:pk}) requires the variability of pharmacy sign-off time, but no direct observations of sign-off are available. The reported model runs therefore use $\kappa = 1.55$, the coefficient of variation calculated from 51 iQemo template-build records as the most complete pharmacy-related field in the tracker. This is a temporary substitution, not a validated estimate of approval-time variability.

The iQemo measure records the time from template set-up to authorisation. Its median is 16 days, and much of that interval is likely to be time spent waiting for the template to be processed. The required approval-time measure, by contrast, concerns the duration of the signing activity itself. The proxy therefore measures a combined waiting-and-signing interval rather than signing time alone.

The direction of the error cannot be established. The observed interval combines waiting and signing time, and adding waiting can change both the mean and the variance. Consequently, the coefficient of variation of the combined interval may be lower than, similar to, or higher than the coefficient of variation of signing alone.

This depends on the relative variability of waiting and signing time, and on whether they move together. Neither is identified in the records. Independence cannot be tested and should not be assumed.

The central problem is identification. The 16-day iQemo interval combines waiting and signing, but the available data do not separate those components. Its coefficient of variation therefore does not identify the coefficient of variation of signing time. The true signing-time $\kappa$ could be lower than, close to, or higher than 1.55. The proxy consequently places little constraint on $\kappa$.

The value is retained solely to allow equation~(\ref{eq:pk}) to be evaluated while direct pharmacy sign-off data are unavailable. It is not treated as an estimate of signing-time variability, and no substantive conclusion in this thesis depends on $\kappa = 1.55$ being accurate.

The observed proxy may not vary in the same way as sign-off time. Direct sign-off observations would be needed to test whether the two are comparable. Pharmacy has therefore been asked for its quickest, usual, and longest hands-on sign-off times. These responses will define plausible bounds rather than provide an exact estimate, and the model will test gate waiting time across those bounds.

Pharmacy's three responses give a minimum-to-maximum range. However, equation~(\ref{eq:pk}) requires the coefficient of variation, $\kappa$, which depends on both the mean and the standard deviation. A range measures the distance between the extremes, not the typical spread of observations around the mean. the range must be converted into an approximate standard deviation before $\kappa$ can be estimated.

Hartley's $d_2$ provides a standard way to convert the range of a normal sample into an estimated standard deviation. Dividing by four is a practical approximation when the range represents roughly 25 observations. Approval times are unlikely to be normally distributed, so this is a convenience assumption rather than a claim about their true distribution. It is used because the available evidence consists only of three broad response from the LTHTR.

Chapter~\ref{ch:results} currently uses $\kappa = 1.55$, because the pharmacy response was not available when the model runs were produced. Once the response is received, the model will be rerun at both ends of the plausible $\kappa$ interval. If the capacity conclusion does not change, the substitution has not affected the result. If it does change, this will show that the conclusion is sensitive to variation in an approval process that the service has not yet measured.

The conversion has three main limitations. First, if approval times are right-skewed, a normal-based conversion may be misleading. Second, the number of observations behind the recalled quickest and longest times is unknown. Third, in a skewed distribution, the range may not be symmetric around the mean. A fitted lognormal or gamma distribution would be preferable, but this would require actual sign-off observations rather than three recalled bands.

A separate unresolved input is $A$, approver availability per cycle (capacity). It is held fixed because the trust has not supplied the required information, and this assumption is documented in Section~\ref{sec:parameters}.

DEAS and GMO committees release studies in batches, which might be conflicted with the model's arrival assumption. However, approval occurs after a long and variable set-up period. That period disperses studies from each committee batch across later months, so much of the original batch pattern is lost before studies reach the approval gate.

This reasoning depends on approval occurring near the end of set-up. It also does not mean that the 95-day duration and the $\kappa$ estimate should be treated as parameters from the same distribution: they are drawn from different samples and should remain separate.
